% This file is part of Oaklisp.
%
% This program is free software; you can redistribute it and/or modify
% it under the terms of the GNU General Public License as published by
% the Free Software Foundation; either version 2 of the License, or
% (at your option) any later version.
%
% This program is distributed in the hope that it will be useful,
% but WITHOUT ANY WARRANTY; without even the implied warranty of
% MERCHANTABILITY or FITNESS FOR A PARTICULAR PURPOSE.  See the
% GNU General Public License for more details.
%
% The GNU GPL is available at http://www.gnu.org/licenses/gpl.html
% or from the Free Software Foundation, 59 Temple Place - Suite 330,
% Boston, MA 02111-1307, USA


\chapter{Numbers} \label{numbers}

\index{\texttt{number}} \index{\texttt{real}} \index{\texttt{complex}}
\index{\texttt{rational}} \index{\texttt{float}} \index{\texttt{integer}}
\index{\texttt{fraction}} \index{\texttt{fixnum}} \index{\texttt{bignum}}

\begin{figure}[h]
\centerline{\input{graphic-numhier}}
\caption{The numeric type hierarchy.  Abstract types are in plain face
and instantiable ones in bold.  Floating point numbers are not
implemented.} \label{fig:numhier}
\end{figure}

\section{Arithmetic}

\op{+}{\dt numbers}
\op{1+}{n}
\op{-}{n1 n2 \dt numbers}
\op{-}{n}
\op{*}{\dt numbers}
\op{/}{n1 n2}
\op{quotient}{n1 n2}
\op{modulo}{n1 n2}
\op{abs}{n1}
\op{max}{n \dt ns}
\op{min}{n \dt ns}
\op{expt}{n1 n2}
\op{gcd}{\dt ns}
\doc{The non-negative greatest common divisor of \emph{ns}, or zero
when there are none.}
\op{lcm}{\dt ns}
\doc{The non-negative least common multiple of \emph{ns}, or one when
there are none.}


\section{Comparison}

\op{=}{n1 n2 \dt ns}
\op{<}{n1 n2 \dt ns}
\op{>}{n1 n2 \dt ns}
\op{<=}{n1 n2 \dt ns}
\op{>=}{n1 n2 \dt ns}
\doc{These take two or more arguments and are true when each adjacent
pair is in the given relation, so \texttt{(< a b c)} means
\texttt{(and (< a b) (< b c))}.  A call with two arguments is open
coded; a longer direct call is open coded as a chain of two-argument
comparisons, and the operations can also be called as values, for
instance through \df{apply}, with any number of arguments from two
up.  Methods added to \df{=} or \df{<} for new types take exactly two
arguments.}

\op{{\protect\bang}=}{n1 n2}
\doc{True when \emph{n1} and \emph{n2} are not equal.  Takes exactly
two arguments.}


\section{Conversion}

\op{string->number}{string $[$radix$]$}
\doc{Parses \emph{string} as a number in \emph{radix}, which defaults
to 10, and returns it, or \df{\#f} if \emph{string} is not a number in
that radix.  The syntax accepted is an optional radix prefix
(\texttt{\#b}, \texttt{\#o}, \texttt{\#d}, \texttt{\#x}), an optional
exactness prefix (\texttt{\#e} or \texttt{\#i}, which is ignored since
all numbers are exact), an optional sign, and an integer, a ratio such
as \texttt{-3/4}, or a decimal such as \texttt{1.5}, read exactly as
\texttt{3/2}, as the reader reads it.}


\section{Predicates}

\pr{zero?}{n}
\pr{negative?}{n}
\pr{positive?}{n}
\pr{even?}{n}
\pr{odd?}{n}


\section{Rounding}

These operations should work on any subtype of \df{real}.

\op{floor}{x}
\doc{Returns the largest integer less than or equal to \emph{x}.}

\op{ceiling}{x}
\doc{Returns the smallest integer greater than or equal to \emph{x}.}

\op{truncate}{x}
\doc{Could be defined \texttt{(if (negative?\ x) (ceiling x) (floor x))}.}

\op{round}{x}
\doc{Returns nearest integer to \emph{x}.  Ties are broken by rounding
to an even number.}

\section{Bitwise Logical Operations}

These operations are only defined for integers.

\op{ash-left}{i amount}
\op{ash-right}{i amount}
\op{rot-left}{i amount}
\op{rot-right}{i amount}
\op{bit-not}{i}
\op{bit-and}{i1 i2}
\op{bit-or}{i1 i2}
\op{bit-nor}{i1 i2}
\op{bit-xor}{i1 i2}
\op{bit-nand}{i1 i2}
\op{bit-andca}{i1 i2}
\op{bit-equiv}{i1 i2}


\section{Accessing Components}

\op{numerator}{rational}
\op{denominator}{rational}

\op{real-part}{number}
\op{imag-part}{number}
